% titel: Blatt 1, A5: Wachstumsfaktoren über zwei Jahre

% Teil I: Werkzeuge, kein Lösungsweg.

% Grundidee: Der mittlere Wachstumsfaktor ist der "geometrische
% Mittelwert" sqrt(ab); er liegt nie über dem arithmetischen Mittel. Warum?
% Ein Quadrat (Differenz zweier Wurzeln) ist nie negativ.

Zwei aufeinanderfolgende Wachstumsfaktoren $a$ und $b$ (z.\,B. ``Faktor
im 1. Jahr``, ``Faktor im 2. Jahr``) wirken \emph{multipliziert}: Über
beide Jahre ist der Gesamtfaktor $ab$. Die Frage ``welcher Faktor wäre
dasselbe Ergebnis in jedem der beiden Jahre?`` führt auf die Zahl $x$ mit
$x^2=ab$, also $x=\sqrt{ab}$ --- den \emph{geometrischen Mittelfaktor}.

\begin{werkzeugkasten}[Mittlerer Wachstumsfaktor]
Wächst eine Population im ersten Jahr um $a>0$, im zweiten um $b>0$,
so ist der Gesamtfaktor $ab$. Der mittlere Wachstumsfaktor $\sqrt{ab}$
hätte in beiden Jahren denselben Gesamtfaktor $(\sqrt{ab})^2=ab$.
Es gilt die zentrale Ungleichung
\begin{equation}
  \sqrt{ab}\le \frac{a+b}2 \qquad (a,b>0),
\end{equation}
der \emph{geometrische} Mittelwert liegt nie über dem
\emph{arithmetischen}.
\end{werkzeugkasten}

\begin{werkzeugkasten}[Warum gilt $\sqrt{ab}\le\frac{a+b}2$? --- der Kernbeweis]
Man startet mit einer trivialen, aber m\"achtigen Tatsache: Ein Quadrat
ist nie negativ.
\begin{equation}
  (\sqrt a-\sqrt b)^2 \ge 0 .
\end{equation}
Ausmultiplizieren: $a-2\sqrt{ab}+b\ge0$, also $a+b\ge2\sqrt{ab}$, geteilt
durch $2$: $\frac{a+b}2\ge\sqrt{ab}$. \textbf{Der Kniff:} man führt die
gewünschte Wurzel $2\sqrt{ab}$ als \emph{Mischterm} eines Quadrats ein,
das man kontrollieren kann. Gleichheit genau dann, wenn
$\sqrt a=\sqrt b$, d.h. $a=b$.
\end{werkzeugkasten}

\begin{achtung}[Darauf achten]
\begin{itemize}
  \item Voraussetzung $a,b>0$ beachten, damit $\sqrt{ab}$ definiert ist.
  \item Die Ungleichung folgt aus dem nichtnegativen Quadrat, nicht aus
        ``Beispielen`` --- Beispiele bestätigen, beweisen aber nicht
        \emph{alte} Behauptung.
  \item Gleichheit gilt genau für $a=b$; schon ein Gegenbeispiel
        ($a\neq b$) zeigt die echte Ungleichung.
\end{itemize}
\end{achtung}
